%% test-japanese.tex
%% =================
%% Typesetting test for the Japanese classes: subcaptions in Japanese,
%% and the space between the label and the text (.5zw).
%%
%%   lualatex test-japanese.tex   (twice)
%%   platex  -jobname=test-japanese "\def\CLS{jsarticle}\input test-japanese"
%%   uplatex -jobname=test-japanese "\def\CLS{ujarticle}\input test-japanese"
%%   (then dvipdfmx test-japanese)
%% The class is ltjsarticle by default; with pTeX the class option
%% dvipdfmx is given.  With \def\SCLlatejapanese{}, the class is
%% article and LuaTeX-ja is loaded after subcaptionlite (LuaLaTeX only):
%%   lualatex -jobname=test-japanese-late \
%%     "\def\SCLlatejapanese{}\input test-japanese"
%%
%% What to check:
%%   * SCL-RESULT OK, and no error or warning.
%%   * In the PDF, the subcaptions are centered under the images, with
%%     a half-width space after (a), (b), ...
\ifdefined\SCLlatejapanese\def\CLS{article}\fi
\ifdefined\CLS\else\def\CLS{ltjsarticle}\fi
\ifdefined\ptexversion\def\CLSOPT{dvipdfmx}\else\def\CLSOPT{}\fi
\expandafter\documentclass\expandafter[\CLSOPT]{\CLS}
%% test-common.tex
%% ===============
%% Helpers shared by the test files; \input this file right after
%% \documentclass, before \usepackage{subcaptionlite}.
%%
%% \SCLlabel{<want>}
%%   Check the label that \label would write here (\@currentlabel), e.g.
%%   "1(a)" for a subfigure or "1" for a figure.
%% \SCLcounter{<counter>}{<value>}
%%   Check the value of <counter> here.
%% \SCLref{<label>}{<want>}
%%   Check the number of <label> that \ref gives (needs a second run).
%% \SCLnoref{<label>}
%%   Check that <label> is not defined (needs a second run).
%% \SCLwd{<want>}{<material>}
%%   Typeset <material> in an \hbox, check its width (<want> is a
%%   dimension) and put the box here.
%% \SCLsep{<kind>}
%%   Check the width of \subcaption@sep in \subcaption@font: <kind> is
%%   "kanji" (half the width of a kanji) or "latin" (.5em).
%% \SCLhref{<want>}
%%   With hyperref, check the anchor of the last \refstepcounter
%%   (\@currentHref); without hyperref, do nothing.
%% \SCLcreftype{<type>}
%%   With cleveref, check the type of the last \refstepcounter for
%%   \cref; without cleveref, do nothing.
%% \@makecaption of the class is saved when this file is read; at the
%% end of \begin{document}, it is checked to be unchanged.
%%
%% Each check logs "SCL-<kind> OK: ..." or "SCL-<kind> NG: ...".  At
%% \end{document}, "SCL-RESULT OK (<n> checks)" or "SCL-RESULT NG" is
%% logged.  Run twice, for \SCLref.
\makeatletter
\newcount\SCL@nchecks
\newif\ifSCL@ng
\def\SCL@ok#1#2{\global\advance\SCL@nchecks\@ne\typeout{SCL-#1 OK: #2}}
\def\SCL@ng#1#2{\global\advance\SCL@nchecks\@ne\global\SCL@ngtrue
  \typeout{SCL-#1 NG: #2}}
%% \SCL@cmp{<kind>}{<got>}{<want>}: compare two strings
\def\SCL@cmp#1#2#3{%
  \edef\SCL@got{\detokenize\expandafter{#2}}%
  \edef\SCL@want{\detokenize{#3}}%
  \ifx\SCL@got\SCL@want
    \SCL@ok{#1}{\SCL@got}%
  \else
    \SCL@ng{#1}{got \SCL@got, want \SCL@want}%
  \fi}
%% \SCL@near{<kind>}{<got>}{<want>}: compare two dimensions (0.01pt)
\def\SCL@near#1#2#3{%
  \begingroup
    \dimen@=\dimexpr#2-(#3)\relax
    \ifdim\dimen@<\z@ \dimen@=-\dimen@ \fi
    \dimen2=#2\relax \dimen4=#3\relax
    \ifdim\dimen@<0.01pt
      \SCL@ok{#1}{\the\dimen2}%
    \else
      \SCL@ng{#1}{got \the\dimen2, want \the\dimen4}%
    \fi
  \endgroup}

\newcommand\SCLlabel[1]{%
  \protected@edef\SCL@tmp{\@currentlabel}%
  \SCL@cmp{LABEL}{\SCL@tmp}{#1}}
\newcommand\SCLcounter[2]{%
  \edef\SCL@tmp{\number\value{#1}}%
  \SCL@cmp{COUNTER}{\SCL@tmp}{#2}}
\newcommand\SCLref[2]{%
  \@ifundefined{r@#1}%
    {\SCL@ng{REF}{#1 undefined}}%
    {\edef\SCL@tmp{\expandafter\expandafter\expandafter\@car
       \csname r@#1\endcsname\@nil}%
     \SCL@cmp{REF}{\SCL@tmp}{#2}}}
\newcommand\SCLnoref[1]{%
  \@ifundefined{r@#1}{\SCL@ok{NOREF}{#1}}{\SCL@ng{NOREF}{#1 defined}}}
\newcommand\SCLwd[2]{%
  \setbox\z@\hbox{#2}%
  \SCL@near{WIDTH}{\wd\z@}{#1}%
  \box\z@}
\newcommand\SCLsep[1]{%
  \begingroup
    \subcaption@font
    \setbox\z@\hbox{\subcaption@sep}%
    \edef\SCL@tmp{#1}\def\SCL@kanji{kanji}%
    \ifx\SCL@tmp\SCL@kanji
      \setbox\tw@\hbox{漢}%
      \dimen@=.5\wd\tw@
    \else
      \dimen@=.5em
    \fi
    \SCL@near{SEP}{\wd\z@}{\dimen@}%
  \endgroup}
\newcommand\SCLhref[1]{%
  \ifdefined\hyper@anchor
    \SCL@cmp{HREF}{\@currentHref}{#1}%
  \fi}
\def\SCL@creftype[#1]#2\@nil{#1}
\newcommand\SCLcreftype[1]{%
  \ifdefined\cref@currentlabel
    \edef\SCL@tmp{\expandafter\SCL@creftype\cref@currentlabel\@nil}%
    \SCL@cmp{CREF}{\SCL@tmp}{#1}%
  \fi}
\let\SCL@makecaption\@makecaption
\AddToHook{begindocument/end}{%
  \ifx\@makecaption\SCL@makecaption
    \SCL@ok{MAKECAPTION}{unchanged}%
  \else
    \SCL@ng{MAKECAPTION}{changed}%
  \fi}

\AddToHook{enddocument/afterlastpage}{%
  \ifnum\SCL@nchecks=\z@ \SCL@ng{RESULT}{no checks}\fi
  \ifSCL@ng
    \typeout{SCL-RESULT NG}%
  \else
    \typeout{SCL-RESULT OK (\the\SCL@nchecks\space checks)}%
  \fi}

%% \SCLimage{<width>}{<height>}: a framed placeholder
\newcommand\SCLimage[2]{\fbox{\rule{0pt}{#2}\hspace{#1}}}
\makeatother
\endinput

\usepackage{subcaptionlite}
\ifdefined\SCLlatejapanese\usepackage{luatexja}\fi
\begin{document}
\section*{副図のテスト}

\SCLsep{kanji}

図~\ref{fig:ja}の左が\ref{fig:ja-a}、右が\ref{fig:ja-b}です。
\SCLref{fig:ja}{1}\SCLref{fig:ja-a}{1(a)}\SCLref{fig:ja-b}{1(b)}

\begin{figure}[htbp]
  \centering
  \begin{subfigure}{.45\linewidth}
    \centering\SCLimage{3cm}{2cm}
    \caption{左の図}\SCLlabel{1(a)}\label{fig:ja-a}
  \end{subfigure}\hfill
  \begin{subfigure}{.45\linewidth}
    \centering\SCLimage{3cm}{2cm}
    \caption{「かぎ括弧」で始まる副図の見出し}\SCLlabel{1(b)}\label{fig:ja-b}
  \end{subfigure}
  \caption{日本語の副図}\SCLlabel{1}\label{fig:ja}
\end{figure}
\SCLcounter{figure}{1}

\begin{figure}[htbp]
  \centering
  \caption{上に置いた図の見出し}\SCLlabel{2}
  \subcaptionbox{副図の見出しが幅より長い場合は折り返されます\SCLlabel{2(a)}}[.3\linewidth]{%
    \SCLimage{3cm}{2cm}}\hfill
  \subcaptionbox{自然幅\SCLlabel{2(b)}}{\SCLimage{3cm}{2cm}}
\end{figure}
\SCLcounter{figure}{2}
\end{document}
